%=====================================================
% CV - Abdoul Hamid Coulibaly (Modern UI/UX Design)
% Generated by scripts/generate-resume.ts — DO NOT EDIT MANUALLY
%=====================================================

\documentclass[10pt,a4paper]{article}

\usepackage[empty]{fullpage}
\usepackage{titlesec}
\usepackage{enumitem}
\usepackage[hidelinks]{hyperref}
\usepackage{fancyhdr}
\usepackage[T1]{fontenc}
\usepackage[utf8]{inputenc}
\usepackage{lmodern}
\usepackage{geometry}
\usepackage{tikz}
\usepackage{xcolor}
\usepackage{tcolorbox}
\usepackage{fontawesome5}
\usepackage{tabularx}
\tcbuselibrary{skins}

% Modern color palette
\definecolor{primary}{RGB}{37, 99, 235}
\definecolor{secondary}{RGB}{71, 85, 105}
\definecolor{darktext}{RGB}{30, 41, 59}
\definecolor{lighttext}{RGB}{55, 65, 81}
\definecolor{bglight}{RGB}{248, 250, 252}
\definecolor{border}{RGB}{226, 232, 240}

% Timeline environments
\newtcolorbox{currentjob}{%
  blanker, left=14pt,
  borderline west={3pt}{0pt}{primary},
  before skip=10pt, after skip=14pt,
  overlay={\fill[primary] ([xshift=1.5pt,yshift=-4pt]frame.north west) circle (6pt);
           \fill[white] ([xshift=1.5pt,yshift=-4pt]frame.north west) circle (3pt);}
}

\newtcolorbox{pastjob}{%
  blanker, left=14pt,
  borderline west={2pt}{0pt}{secondary!50},
  before skip=10pt, after skip=14pt,
  overlay={\fill[secondary!50] ([xshift=1pt,yshift=-3pt]frame.north west) circle (5pt);
           \fill[white] ([xshift=1pt,yshift=-3pt]frame.north west) circle (2.5pt);}
}

% Skill tag
\newtcbox{\skilltag}{on line, arc=4pt, colback=bglight, colframe=bglight,
  boxrule=0pt, boxsep=0pt, left=5pt, right=5pt, top=2pt, bottom=2pt,
  fontupper=\footnotesize\color{darktext}}

\newtcbox{\primarytag}{on line, arc=4pt, colback=primary!10, colframe=primary!10,
  boxrule=0pt, boxsep=0pt, left=5pt, right=5pt, top=2pt, bottom=2pt,
  fontupper=\footnotesize\color{primary}\bfseries}

\geometry{left=1.5cm, right=1.5cm, top=1cm, bottom=1cm, footskip=5pt}

\pagestyle{fancy}
\fancyhf{}
\renewcommand{\headrulewidth}{0pt}

\setlength{\parindent}{0pt}
\setlist[itemize]{itemsep=3pt, topsep=4pt, parsep=1pt, leftmargin=1em, label=\textcolor{primary}{\textbullet}}

\titleformat{\section}
{\large\bfseries\color{darktext}}{}{0em}{}[\vspace{-6pt}\textcolor{border}{\rule{\linewidth}{1pt}}]
\titlespacing*{\section}{0pt}{12pt}{8pt}

%=====================================================
\begin{document}

%================ HEADER ==================
\begin{center}
{\fontsize{22}{26}\selectfont\textbf{\textcolor{darktext}{Abdoul Hamid COULIBALY}}}\\[3pt]
{\large\textcolor{primary}{\textbf{Backend Engineer}} \textcolor{lighttext}{---} \textcolor{secondary}{Syst\`emes Data \& Cloud}}\\[8pt]
\textcolor{secondary}{%
\faMapMarker*\ Rennes, France \quad
\faPhone*\ +33 7 49 10 66 71 \quad
\faEnvelope\ \href{mailto:pro@dimahc.dev}{pro@dimahc.dev} \quad
\faGlobe\ \href{https://dimahc.dev}{dimahc.dev} \quad
\faLinkedin\ \href{https://linkedin.com/in/abdoul-hamid-coulibaly}{LinkedIn}
}
\end{center}

\vspace{4pt}

\begin{tcolorbox}[colback=bglight, colframe=bglight, arc=4pt, boxrule=0pt, left=10pt, right=10pt, top=6pt, bottom=6pt]
\small\textcolor{darktext}{Ing\'enieur backend Go et Python avec 4 ans d'exp\'erience sur des plateformes de donn\'ees distribu\'ees. Ingestion temps r\'eel, microservices, architectures event-driven. J'ai travaill\'e aussi bien sur la conception de pipelines que sur leur d\'eploiement et leur tenue en production~: astreinte, r\'esolution d'incidents, observabilit\'e.}
\end{tcolorbox}

%=====================================================
\section*{Exp\'erience Professionnelle}

\begin{currentjob}
\textbf{\large\textcolor{darktext}{Ingénieur Backend -- Data Supply Chain}} \hfill \colorbox{primary}{\textcolor{white}{\small\bfseries\textsf{\ Novembre 2023 - Présent\ }}}\\[2pt]
\textcolor{primary}{\textbf{Univers (QoS Energy)}} \textcolor{lighttext}{| La Chapelle-sur-Erdre, France}

\vspace{4pt}

\textbf{\textcolor{secondary}{Architecture \& Ingestion de Données}}
\begin{itemize}
\item Développement et refonte de microservices Go et Python pour l'ingestion de millions de métriques par minute, dans un environnement Kubernetes géré en GitOps.
\item Mise en place et extension des architectures event-driven (RabbitMQ, NATS) pour l'ingestion et le traitement asynchrone.
\item Conception de pipelines résilients face aux aléas réseau et à l'instabilité des sources : gestion d'erreurs distribuées, retry, backfilling et rechargement manuel.
\end{itemize}

\vspace{6pt}
\textbf{\textcolor{secondary}{Intégration \& Connecteurs}}
\begin{itemize}
\item Développement de connecteurs pour sources hétérogènes : SCADA, REST, FTP, bases SQL.
\item Participation active aux décisions d'architecture : évaluation de technologies (NATS, Redis), modélisation des données et gestion du rate limiting des APIs sources.
\item Réécriture de services Go en Python pour s'aligner sur les contraintes stack de Bazefield.
\end{itemize}

\vspace{6pt}
\textbf{\textcolor{secondary}{Scalabilité \& Résilience}}
\begin{itemize}
\item Instrumentation de services avec Prometheus et mise en place de dashboards Grafana pour le suivi opérationnel.
\item Astreinte opérationnelle (1 sem./6) : résolution d'incidents critiques et support L3 pour les opérations techniques clients.
\end{itemize}

\vspace{3pt}
\primarytag{Go} \skilltag{Python} \skilltag{Ruby} \skilltag{RabbitMQ} \skilltag{NATS} \skilltag{PostgreSQL} \skilltag{Kubernetes} \skilltag{Azure} \skilltag{GitOps} \skilltag{Prometheus} \skilltag{Grafana}
\end{currentjob}

\begin{pastjob}
\textbf{\large\textcolor{darktext}{Ingénieur Logiciel}} \hfill {\small\textsf{Octobre 2021 - Octobre 2023}}\\[2pt]
\textcolor{secondary}{\textbf{NumoData (ex-EXFO Solutions)}} \textcolor{lighttext}{| Rennes, France}

\vspace{4pt}

\textbf{\textcolor{secondary}{Développement Backend}}
\begin{itemize}
\item Développement de services backend Go pour une plateforme de supervision réseau traitant des millions de métriques par minute.
\item Conception de requêtes analytiques (SQL, Druid) : agrégations, requêtes multidimensionnelles, optimisation via indexation et cache.
\end{itemize}

\vspace{6pt}
\textbf{\textcolor{secondary}{Smart Monitoring d'Usage --- Ownership Technique}}
\begin{itemize}
\item Conception et développement end-to-end d'un service d'analyse d'usage des dashboards : métriques les plus consultées, dimensions les plus volumineuses, détection des données sous-utilisées.
\item Exposition des résultats via API REST, instrumentation Prometheus et dashboards Grafana, exploitation pour l'optimisation du cache et la réduction des coûts de stockage.
\end{itemize}

\vspace{6pt}
\textbf{\textcolor{secondary}{Performance \& Infrastructure}}
\begin{itemize}
\item Conduite de tests de charge et analyse de corrélation volumétrie/ressources (CPU, RAM, latence) pour le dimensionnement client.
\item Développement d'un outil interne de génération automatique de configuration Kubernetes (Helm) à partir des résultats de dimensionnement client.
\end{itemize}

\vspace{3pt}
\primarytag{Go} \skilltag{Python} \skilltag{Kafka} \skilltag{Apache Druid} \skilltag{AWS} \skilltag{Prometheus} \skilltag{Grafana} \skilltag{Helm} \skilltag{Kubernetes}
\end{pastjob}

\begin{pastjob}
\textbf{\large\textcolor{darktext}{Développeur Web -- Stage}} \hfill {\small\textsf{Mai 2021 - Août 2021}}\\[2pt]
\textcolor{secondary}{\textbf{Resalocal}} \textcolor{lighttext}{| Talloires-Montmin, France}

\vspace{4pt}

\begin{itemize}
\item Développement d'un pipeline d'ingestion depuis la base touristique Apidae vers le PIM Akeneo : import de POIs, événements et activités.
\item Synchronisation bidirectionnelle : restitution des données enrichies vers Apidae après traitement.
\item Gestion des médias associés aux POIs via l'API Cloudinary : upload et transformation d'images.
\item Gestion des contraintes d'ingestion : pagination, rate limiting, et préservation des données existantes dans le PIM.
\end{itemize}

\vspace{3pt}
\primarytag{PHP} \skilltag{Laravel} \skilltag{Docker} \skilltag{Kubernetes} \skilltag{Akeneo} \skilltag{Cloudinary} \skilltag{API REST}
\end{pastjob}

%=====================================================
\section*{Comp\'etences Techniques}

\vspace{4pt}

\begin{tabularx}{\textwidth}{@{}l X@{}}
\textcolor{primary}{\faCode}\ \textbf{Langages} & \skilltag{Go} \skilltag{Python} \skilltag{Ruby} \skilltag{TypeScript} \skilltag{SQL} \skilltag{Bash} \\[4pt]
\textcolor{primary}{\faCloud}\ \textbf{Cloud \& DevOps} & \skilltag{AWS} \skilltag{Azure} \skilltag{Docker} \skilltag{Kubernetes} \skilltag{Terraform} \skilltag{Helm} \skilltag{GitLab CI} \\[4pt]
\textcolor{primary}{\faDatabase}\ \textbf{Data \& Messaging} & \skilltag{PostgreSQL} \skilltag{MySQL} \skilltag{Apache Druid} \skilltag{Kafka} \skilltag{RabbitMQ} \skilltag{NATS} \\[4pt]
\textcolor{primary}{\faChartLine}\ \textbf{Observabilit\'e} & \skilltag{Prometheus} \skilltag{Grafana} \skilltag{OpenTelemetry}
\end{tabularx}

%=====================================================
\section*{Formation}

\vspace{4pt}

\begin{tabularx}{\textwidth}{@{}X r@{}}
\textbf{Master MIAGE -- DABI (Décisionnel, Apprentissage et Big Data)} \textcolor{lighttext}{| Université de Rennes 1} & \textcolor{lighttext}{2021 -- 2023} \\
\multicolumn{2}{@{}p{\textwidth}@{}}{\textcolor{lighttext}{\small Entrepôts de données, fouille de données, Cloud \& Big Data, intelligence artificielle, intégration d'applications}} \\[4pt]
\textbf{Licence Informatique -- MIAGE} \textcolor{lighttext}{| Université de Rennes 1} & \textcolor{lighttext}{2020 -- 2021} \\
\multicolumn{2}{@{}p{\textwidth}@{}}{\textcolor{lighttext}{\small Double compétence informatique/gestion, bases de données, programmation web, modélisation objet, analyse de données}} \\[4pt]
\textbf{Diplôme de Technicien Supérieur en Informatique} \textcolor{lighttext}{| INP-HB} & \textcolor{lighttext}{2016 -- 2019} \\
\multicolumn{2}{@{}p{\textwidth}@{}}{\textcolor{lighttext}{\small Algorithmique, réseaux, développement web, administration systèmes, gestion de projets}} \\[4pt]
\textbf{Formation Anglais Général et Professionnel} \textcolor{lighttext}{| Wall Street English} & \textcolor{lighttext}{2024} \\
\multicolumn{2}{@{}p{\textwidth}@{}}{\textcolor{lighttext}{\small Communication orale et écrite, vocabulaire technique IT, présentations professionnelles, réunions en anglais}}
\end{tabularx}

%=====================================================
\vspace{6pt}
\section*{Langues}

\vspace{4pt}

\textbf{Fran\c{c}ais} \textcolor{lighttext}{(natif)} \quad | \quad
\textbf{Anglais} \textcolor{lighttext}{(courant)} \quad | \quad
\textbf{Bambara} \textcolor{lighttext}{(natif)}

%=====================================================
\end{document}
